\subsection{幺模群 SL(n, A)}

设 \(\mathbf{M}_n(\mathbf{A})\) 表示 \(n\) 阶方阵（ \(\mathbf{A}\) 上的）构成的环。考虑映射 \(\det: \mathbf{M}_n(\mathbf{A}) \to \mathbf{A}\) 。群 \(\mathbf{GL}(n, \mathbf{A})\) （ \(\mathbf{M}_n(\mathbf{A})\) 的可逆元素组成的群，同构于 A-模 \(\mathbf{A}^n\) 的自同构群（II，§ 10，第 7 号））正是在该映射下乘法群 \(\mathbf{A}^*\) （ \(\mathbf{A}\) 的可逆元素组成的群）的逆像（第 3 号，命题 5）。另一方面注意，映射 \(\det: \mathbf{GL}(n, \mathbf{A}) \to \mathbf{A}^*\) 是一个群同态（第 3 号，命题 5）。

映射 \(\operatorname{det}:\mathbf{M}_n(\mathbf{A})\to \mathbf{A}\) 此外是满射（因此同态 \(\operatorname{det}:\mathbf{GL}(n,\mathbf{A})\to \mathbf{A}^*)\) 也是满射；因为对所有 \(\lambda \in \mathbf{A}\) ，

\[
\det (\operatorname{diag} (\lambda , 1, \dots , 1)) = \lambda
\]

这由第 6 号的公式 (32) 得出。

满射同态 \(\operatorname{det}:\mathbf{GL}(n,\mathbf{A})\to \mathbf{A}^*\) 的核是 \(\mathbf{GL}(n,\mathbf{A})\) 的一个正规子群，它由幺模矩阵组成；它记作 \(\mathbf{SL}_n(\mathbf{A})\) 或 \(\mathbf{SL}(n,\mathbf{A})\) ，常称为 \(n\) 阶方阵（ \(\mathbf{A}\) 上的）的幺模群或特殊线性群。

在本号中，我们将考察 A 为域的情形。回忆对于 \(1 \leqslant i \leqslant n\) , \(1 \leqslant j \leqslant n\) , \(E_{ij}\) 表示 \(n\) 阶方阵，其所有元素均为零，除了指标 \(i\) 的行与指标 \(j\) 的列交处等于 1 的那个元素；以 \(I_n\) 表示 \(n\) 阶单位矩阵，我们记 \(B_{ij}(\lambda) = I_n + \lambda E_{ij}\) （对每一对不同的有序指标 \(i, j\) 以及所有 \(\lambda \in A\) ；II，§ 10，第 13 号）。

命题 17. 设 K 是一个交换域。幺模群 \(\mathbf{SL}(n, \mathbf{K})\) 由矩阵 \(B_{ij}(\lambda)\) （ \(i \neq j\) ， \(\lambda \in \mathbf{K}\) ）生成。

由 II，§ 10，第 13 号，命题 14，我们知道 \(\mathbf{GL}(n,\mathbf{K})\) 中的每个矩阵都是形如 \(B_{ij}(\lambda)\) 的矩阵与一个形如 \(\operatorname{diag}(1,1,\ldots,1,\alpha)\) （其中 \(\alpha\in\mathbf{K}^{*}\) ）的矩阵的乘积。现在显然有 \(\det(B_{ij}(\lambda))=1\) 及 \(\det(\operatorname{diag}(1,\ldots,1,\alpha))=\alpha\) （第 6 号，例 2）；由此得出该命题。

推论. 群 \(\mathbf{SL}(n, \mathbf{K})\) 是 \(\mathbf{GL}(n, \mathbf{K})\) 的换位子群，除非 \(n = 2\) 且 \(\mathbf{K}\) 是含 2 个元素的域。

由于 \(\mathbf{SL}(n, \mathbf{K})\) 是从 \(\mathbf{GL}(n, \mathbf{K})\) 到交换群 \(\mathbf{K}^*\) 的 det 同态的核， \(\mathbf{SL}(n, \mathbf{K})\) 包含换位子群 \(\Gamma\) （ \(\mathbf{GL}(n, \mathbf{K})\) 的）（I，§ 6，第 2 号）。为证明 \(\mathbf{SL}(n, \mathbf{K}) = \Gamma\) ，根据命题 17，只需证明对所有 \(\lambda \in \mathbf{K}^*\) ， \(B_{ij}(\lambda)\) 属于 \(\Gamma\) 。现在， \(B_{ij}(\lambda)\) 是 \(B_{ij}(1)\) 在 \(\mathbf{GL}(n, \mathbf{K})\) 中的共轭，因为 \(B_{ij}(\lambda) = Q \cdot B_{ij}(1) \cdot Q^{-1}\) ，其中 \(Q\) 表示典范基（ \(e_i\) ）下自同构 \(v\) （ \(\mathbf{K}^n\) 的）的矩阵，它使得 \(v(e_i) = \lambda e_i\) , \(v(e_k) = e_k\) （ \(k \neq i\) ）。另一方面，设 \(u_{ij}\) （对于 \(i \neq j\) ）为 \(\mathbf{K}^n\) 的自同构，它使得 \(u_{ij}(e_i) = -e_j\) , \(u_{ij}(e_j) = e_i\) , \(u_{ij}(e_k) = e_k\) （ \(k \notin \{i, j\}\) ），且属于 \(\mathbf{SL}(n, \mathbf{K})\) ；于是 \(B_{ji}(\lambda) = U_{ij}B_{ij}(-\lambda)U_{ij}^{-1}\) ，其中 \(U_{ij}\) 是 \(u_{ij}\) 关于典范基的矩阵。类似地，如果 \(1 < i < j\) ，则 \(B_{1j}(\lambda) = U_{1i}B_{ij}(\lambda)U_{1i}^{-1}\) ；最后，对于 \(2 < j\) ， \(B_{12}(\lambda) = U_{2j}B_{1j}(\lambda)U_{2j}^{-1}\) 。这证明所有 \(B_{ij}(\lambda)\) 具有相同的像 \(s\) （在 \(\mathbf{GL}(n, \mathbf{K}) / \Gamma\) 中），并且仍需证明 \(s\) 是单位元。

首先假设 K 包含一个异于 0 和 1 的元素 \(\lambda\) ；那么 \(1 = \lambda + (1 - \lambda)\) ，右端的两个项均 \(\neq 0\) ；关系 \(B_{12}(1) = B_{12}(\lambda)B_{12}(1 - \lambda)\) 表明 \(s^{2} = s\) ，从而 s 是单位元。

现在假设 \(n \geqslant 3\) 。乘积 \(B_{21}(1)B_{31}(1)\) 是自同构 \(u\) （ \(\mathbf{K}^n\) 的）的矩阵，它使得 \(u(e_1) = e_1 + e_2 + e_3\) 与 \(u(e_i) = e_i\) （对 \(i \neq 1\) ）。若 \(S\) 是自同构 \(u'\) （ \(\mathbf{K}^n\) 的）的矩阵，它使得 \(u'(e_2) = e_2 + e_3\) 与 \(u'(e_i) = e_i\) （对 \(i \neq 2\) ），则 \(S.B_{21}(1)B_{31}(1).S^{-1} = B_{21}(1)\) ；我们还推出 \(s^2 = s\) ，这就完成了证明。

注记。(1) \(\mathbf{GL}(2, \mathbf{F}_2) = \mathbf{SL}(2, \mathbf{F}_2)\) ；这是一个阶为 6 的可解群，其换位子群的指数为 2（II，§ 10，习题 14）。

\begin{enumerate}
\def\labelenumi{(\arabic{enumi})}
\setcounter{enumi}{1}
\tightlist
\item
  沿用上述记号，可如同I，§5，第7号，命题9那样证明，对于 \(i < j\) , \(j - i > 1\) , \(u_{ij} = u_{j-1,j} u_{i,j-1} u_{j-1,j}^{-1}\) ；因此群 \(\mathbf{SL}(n, \mathbf{K})\) 由矩阵 \(B_{12}(\lambda)\) 与诸 \(U_{i,i+1}\) （ \(1 \leqslant i \leqslant n - 1\) ）生成。
\end{enumerate}

